% !TEX root = ../nc-quadric-classification.tex
% 2026-11-21 / v1.1
\section{有限 foundation と反 Serre 表示のスタック}
\label{sec:foundation-stack}

この節の有限代数スタックは直ちに構成できる。
それを $\HH$ の全相対モジュライと同一視するためには、末尾に記す
相対存在・開性・降下の証明がさらに必要である。

\subsection{二つの線形写像によるパラメータ}

$P,Q$ を二次元、$W_1,W_2,W_3$ をそれぞれ二、四、六次元のベクトル空間とする。
箙 $1\overset{2}{\longrightarrow}2\overset{2}{\longrightarrow}3$ の道代数を $H$ と書く。
隣接矢空間は $P,Q$ であり、$H_{13}=P\otimes Q$ である。
次元ベクトル $(2,4,6)$ の $H$-表現は
\[
 F:P\otimes W_1\longrightarrow W_2,\qquad
 G:Q\otimes W_2\longrightarrow W_3
\]
の組で与えられる。従ってそのパラメータ空間は
\[
 \mathcal R_{\mathrm{fnd}}
 =\Hom(P\otimes W_1,W_2)\oplus\Hom(Q\otimes W_2,W_3)
 \simeq\mathbb A^{64}.
\]
表現 $M=(W_1,W_2,W_3;F,G)$ に対応する一頂点拡大を $\Lambda_M$ とする。
具体的には頂点を $0,1,2,3$ とし、
\[
 (\Lambda_M)_{0i}=W_i\ (1\le i\le3),\qquad
 (\Lambda_M)_{ij}=H_{ij}\ (1\le i\le j\le3),\qquad
 (\Lambda_M)_{00}=k
\]
と置く。積は $H$ の道の積と表現 $M$ の作用による。
これは二十四次元の有向代数である。
箙の関係式の有限型スタックの一般的構成は
Abdelgadir--Okawa--Ueda \cite[Definition 3.5, Proposition 3.6]{AOU2014} にある。
ここでは全ての合成則を明示し、結合律による方程式も消去できる。

群
\[
 \mathcal G=\GL(W_1)\times\GL(W_2)\times\GL(W_3)
             \times\GL(P)\times\GL(Q)
\]
を、$g=(g_1,g_2,g_3,a,b)$ に対し
\[
 (F,G)\longmapsto
 \bigl(g_2F(a^{-1}\otimes g_1^{-1}),\
       g_3G(b^{-1}\otimes g_2^{-1})\bigr)
\]
で作用させる。中心トーラス
\[
 K=\left\{
 \left(\lambda_1 I,\lambda_2 I,\lambda_3 I,
       \frac{\lambda_2}{\lambda_1}I,
       \frac{\lambda_3}{\lambda_2}I\right)
 \right\}\simeq\mathbb G_m^3
\]
は全空間で自明に作用する。

\begin{proposition}[有限 foundation の商スタック]
\label{prop:foundation-stack}
商
\[
 \mathfrak F^{\mathrm{lin}}=[\mathcal R_{\mathrm{fnd}}/\mathcal G],\qquad
 \mathfrak F=[\mathcal R_{\mathrm{fnd}}/(\mathcal G/K)]
\]
は有限型の滑らかな Artin スタックである。
前者は上の形の有向代数を頂点とともに記述し、後者は中心 $K$ の rigidification である。
$\dim\mathfrak F=3$ である。
全ての $\HH$ 対象の定理~\ref{thm:normal-helix} の foundation は、
その End 代数を通じて $\mathfrak F^{\mathrm{lin}}$ の幾何点を与える。
\end{proposition}

\begin{proof}
$\mathcal R_{\mathrm{fnd}}$ は滑らかな六十四次元アフィン空間であり、
$\dim\mathcal G=4+16+36+4+4=64$ である。
$K$ は滑らかな中心トーラスなので $\mathcal G/K$ も滑らかな代数群である。
従って二つの商は滑らかな Artin スタックであり\cite[Tag 0DLS]{StacksModuli}、
$\dim\mathfrak F=64-(64-3)=3$ である。

頂点を保つ $H$ の自己同型は $P,Q$ 上の基底変更で尽くされる。
したがって $\mathcal G$ の作用は、頂点を保つ $\Lambda_M$ の同型を正確に与える。
対角単元 $\sum_{i=0}^3\lambda_i e_i$ による内部自己同型は、共通倍率で割り
$\lambda_0=1$ とすれば上の $K$ となる。
これは各 foundation 対象の比較同型のスカラーによる変更を、全体の自然スカラーで割ったものに対応する。
従って $K$ の除去は、$E_0$ の標識だけを残して、
$E_1,E_2,E_3$ を保つ比較同型の選択を忘れる rigidification である。
幾何点では、foundation の各項を同型類として保つ Morita 同値の自己同型群が
$\operatorname{Aut}_{\{e_i\}}(\Lambda_M)/K$ となる。
任意の底上では、これらの対象の存在と線束 twist の降下を区別する必要がある。

$\HH$ から得た四項 $(E_0,E_1,E_2,E_3)$ では
\[
 \dim\Hom(E_i,E_j)=
 \begin{pmatrix}
 1&2&4&6\\0&1&2&4\\0&0&1&2\\0&0&0&1
 \end{pmatrix}_{ij}
\]
であり、$\Hom(E_1,E_2)\otimes\Hom(E_2,E_3)\to\Hom(E_1,E_3)$ は同型である。
後者は系~\ref{cor:normal-cubic-defect} の奇数頂点間の積の計算による。
よって $E_1,E_2,E_3$ の End 代数が $H$ となり、
$\Hom(E_0,E_i)$ が上の表現 $M$ を与える。
\end{proof}

\begin{remark}
$\mathfrak F$ は foundation の線形化を除いた有限代数のパラメータスタックである。
任意の点の $\operatorname{Perf}(\Lambda_M)$ に、$\HH$ の heart、
相対 Serre 根、豊富性、AS 正則被覆が存在するとは主張していない。
また foundation を保たない全ての pointed 同値を、この群 $\mathcal G/K$ だけで割ったとも主張しない。
一般に有限代数の導来圏の同値と元の heart の exact 同値は区別する。
本稿の体上では構造対象を保つ導来同値の exact 性を
系~\ref{cor:pointed-derived-exact} で証明しているが、族での底変更には別途確認を要する。
\end{remark}

\subsection{Hochschild 障害の消滅}

\begin{proposition}
\label{prop:foundation-hochschild}
任意の $M\in\mathcal R_{\mathrm{fnd}}$ に対し
\[
 \operatorname{HH}^n(\Lambda_M)=0\quad(n\ge3),\qquad
 \dim\operatorname{HH}^2(\Lambda_M)
 -\dim\operatorname{HH}^1(\Lambda_M)=3.
\]
従ってこれらの代数の通常の結合代数変形について、三次 Hochschild コホモロジーにある障害は消える。
これは相対 $\HH$ 条件の無障害性とは別の主張である。
\end{proposition}

\begin{proof}
$A=\Lambda_M$、$E=\bigoplus_{i=0}^3ke_i\simeq k^4$ とする。
$E$ は分離代数なので、$E$ に相対的な正規化 bar 複体が $\operatorname{HH}^*(A)$ を計算する。
有向性からそのコチェイン項は
\[
 C^n=\bigoplus_{i_0<\cdots<i_n}
 \Hom_k\bigl(A_{i_0i_1}\otimes\cdots\otimes A_{i_{n-1}i_n},
             A_{i_0i_n}\bigr)
\]
であり、$C^0=k^4$、$C^n=0$ ($n>3$) である。
特に
\[
 C^3=\Hom(A_{01}\otimes A_{12}\otimes A_{23},A_{03}).
\]
$C^2$ の $(0,1,3)$ 成分
$\Hom(A_{01}\otimes A_{13},A_{03})$ から $C^3$ への微分を取る。
他の成分を零とすれば、bar 微分の四項のうち残るものは
\[
 (a,b,c)\longmapsto\varphi(a,bc)
\]
だけである。$A_{12}\otimes A_{23}\xrightarrow\sim A_{13}$ なので
この写像は二つの四十八次元空間の同型である。
従って $C^2\to C^3$ は全射で、$\operatorname{HH}^{\ge3}(A)=0$ となる。

各コチェイン項の次元は
\[
 (\dim C^0,\dim C^1,\dim C^2,\dim C^3)
 =(4,80,128,48).
\]
ここで $\dim C^1=2^2+4^2+6^2+2^2+4^2+2^2=80$、
$\dim C^2=16+48+48+16=128$ である。
中心の元は全 $e_i$ と可換なので対角であり、
各 $A_{0i}=W_i\ne0$ と可換であるため四つの対角係数は全て等しい。
従って $\operatorname{HH}^0(A)=Z(A)=k$ である。
Euler 標数を取ると
\[
 1-\dim\operatorname{HH}^1(A)+\dim\operatorname{HH}^2(A)
 =4-80+128-48=4
\]
となり、第二式が従う。
\end{proof}



\subsection{二つの ruling と foundation の内在性}

\begin{proposition}[ruling 対象の内在的認識]
\label{prop:intrinsic-ruling-pair}
$(\mathcal C,\mathcal O)\in\HH$、$t=S^{-1}[2]$ とする。
Gabriel--Krull rank one torsion-free exceptional 対象 $F$ で
\[
 \chi(\mathcal O,t^nF)=(2n+1)(2n+2)\qquad(n\gg0)
\]
を満たすものの同型類は、正確に二つである。
任意の適合表示の二つの次数零生成 MCM 型の tails 像が、この二類を与える。
従ってその非順序対は $(\mathcal C,\mathcal O)$ から内在的に定まる。
\end{proposition}

\begin{proof}
適合表示を一つ選び、二つの MCM 型の tails 像を $X,Y$ と書く。
定理~\ref{thm:normal-helix} の評価列により、次数付き加群 $Z=X,Y$ と
$p(Z)$ は自由加群に埋め込まれる。
$B$ の商斜体へ移せば両者の正の rank の和は二なので、各 rank は一である。
従って $Z(-1)\hookrightarrow B^2$ の非零成分 $Z(-1)\to B$ は単射である。
実際、商斜体へ移した核は零であり、$Z$ は自由加群の部分加群なので
この核はもともと零である。
また $Z$ の非零斉次元は単射 $B(-d)\hookrightarrow Z$ を与える。
tails と Gabriel--Krull 商 $q:\mathcal C\to\mathcal C/\mathcal C_1$ に移すと
\[
 q\pi B(-d)\hookrightarrow q\pi Z\hookrightarrow q\pi B(1)
\]
となり、外側は非零単純対象なので $q\pi Z$ も非零単純である。
$\pi Z\hookrightarrow\pi B(1)$ によって torsion-free 性も従う。
exceptionality は同定理の geometric helix から従い、
Hilbert 級数 $2/(1-z)^3$ から指定の Euler 多項式を得る。

命題~\ref{prop:first-orbit-object} と同じ順序で
$K_0(\mathcal C)$ を $K_0(\mathbb P^1\times\mathbb P^1)$ と同定する。
$[F]$ の Chern 指標を $r+aH_1+bH_2+cH_1H_2$ と書くと
\[
 \chi(\mathcal O,t^nF)
 =r(2n+1)^2+(a+b)(2n+1)+c,\qquad
 \chi(F,F)=r^2+2rc-2ab.
\]
係数比較と例外性から $r=1$、$a+b=1$、$c=0$、$ab=0$ となる。
従って可能な類は $(a,b)=(1,0),(0,1)$ の二つだけである。

同じ $K_0$ 類を持つ二つの rank one torsion-free exceptional 対象 $L,F$ が
この多項式を満たすなら $\chi(L,F)=1$ である。
もし $\Ext^2(L,F)\ne0$ なら Serre 双対性から $F\to t^{-1}L$ が非零であり、
rank one torsion-free 性により単射となる。
十分大きい $n$ で Hom 次元は Euler 特性と一致する。
この点は命題~\ref{prop:first-orbit-object} の証明の自由被覆の議論による。
具体的には任意の $T$ を $\mathcal O$ のシフトの有限直和 $P$ で覆い、
核も同様に $P'$ で覆う。自由シフト間の $\Ext^1$ の消滅から
$\Ext^1(t^{-n}\mathcal O,T)$ は $\Ext^2(t^{-n}\mathcal O,P')$ の部分商であり、
後者は Serre 双対性と負次数切断の消滅で零となる。
$\Ext^2(t^{-n}\mathcal O,T)$ についても $P$ で同じ議論が使える。
従って上の単射は
\[
 (2n+1)(2n+2)\le(2n-1)(2n)
\]
を与えて矛盾する。従って $\Hom(L,F)\ne0$ であり、同じ理由で単射である。
その余核 $T$ は $K_0$ 類が零なので Hilbert 多項式が零である。
ここで有限生成 $B$-加群を $A$-加群と見ると有限自由分解を持つため
Hilbert 関数は最終的に多項式であり、有限次元加群を除いてもその多項式は変わらない。
従って零 Hilbert 多項式を持つ加群は有限次元であり、$T=0$ となる。
よって各 $K_0$ 類に高々一つの同型類しかない。
$X,Y$ は互いに非同型なので二つの類を尽くす。
\end{proof}

\begin{corollary}[指定構造対象から始まる二つの foundation]
\label{cor:two-intrinsic-foundations}
$M=s\mathcal O$ を命題~\ref{prop:first-orbit-object} の内在的対象とする。
定理~\ref{thm:normal-helix} で作る、添字零が $\mathcal O$ である foundation は、
各項の同型まで正確に二つである。
それぞれ最初の三項が $(\mathcal O,X,M)$ または $(\mathcal O,Y,M)$ となり、
第四項はこれらの三項から内在的に決まる。
pointed 自己同値群はこの二元集合に作用する。
\end{corollary}

\begin{proof}
初めの三項の内在性は直前の命題と第一軌道対象の認識による。
一つの選択 $E_1=X$ または $Y$ を固定する。
既存の full foundation の完全性から
\[
 {}^{\perp}\langle\mathcal O,E_1,M\rangle\simeq D^b(\operatorname{mod}k).
\]
この左直交の exceptional 生成元は、第四項の cohomological shift に限られる。
条件 $\operatorname{RHom}(M,E_3)\simeq k^2[0]$ はその shift を一意に固定する。
従って第四項も一意である。
\end{proof}

\begin{remark}[相対的な二重被覆へ残る点]
上の二択は代数閉体上の対象についての結論である。
相対的な有限 étale 二重被覆、あるいは単一の $S_2$ 商による
$\HH$ スタックの構成には、次を別途示す必要がある。
\begin{enumerate}
\item 相対 Serre と Euler 多項式によって定義した $M,X,Y$ の
      fppf 局所的存在、底変更、線束 twist を除く降下。
\item $\HH$ に対応する $\mathcal R_{\mathrm{fnd}}$ の部分関手が
      有限表示の開または局所閉条件となること。
\item ruling の交換で得る新しい foundation が族で代数的に構成され、
      二回の交換と恒等の整合性を持つこと。
\item 相対 AZ 再構成と heart の内在的認識により、
      このgroupoidの全矢が元の $\mathcal O$-marked exact 同値を与えること。
\end{enumerate}
exceptional 対象の $\Ext^1$ と $\Ext^2$ の消滅は、微小拡大上の一意延長の
根拠になるが、上記の有限表示・開性・代数化を単独で証明するものではない。
$\mathfrak F$ の滑らかさや $\operatorname{HH}^3$ の消滅だけから、
$\HH$ の相対条件が全ての変形で保たれるとは結論しない。
\end{remark}

\subsection{反 Serre 環の商と標識に由来する慣性}
\label{subsec:antiserre-stack-candidate}

ここでは $k$ 上の pointed 圏の射を
$(F,u):(\mathcal C,\mathcal O)\to(\mathcal C',\mathcal O')$、
$u:F\mathcal O\xrightarrow\sim\mathcal O'$ とし、
二つの射の間の二射を $u'\eta_{\mathcal O}=u$ を満たす自然同型
$\eta:F\Rightarrow F'$ とする。
以下の自己同型群は、この二groupoidの一射を二射で割った群を意味する。

\begin{lemma}[標識付き圏の二射と次数スカラー]
\label{lem:pointed-degree-scalar}
$(\mathcal C,\mathcal O)\in\HH$ とし、
$R=\bigoplus_{n\ge0}\Hom(\mathcal O,t^n\mathcal O)$ とする。
$\operatorname{Aut}(\operatorname{Id}_{\mathcal C})=k^\times$ であり、
pointed 一射の二自己同型は自明である。
また完全列
\[
 1\longrightarrow\mathbb G_m(k)
 \xrightarrow{\lambda\mapsto\xi_\lambda}
 \operatorname{Aut}_{\mathrm{gr}}(R)
 \longrightarrow
 \operatorname{Aut}_{\mathrm{pt}}(\mathcal C,\mathcal O)
 \longrightarrow1,
 \qquad \xi_\lambda|_{R_n}=\lambda^n\operatorname{id}
\]
がある。従って graded algebra の商をそのまま用いると、
pointed 圏の自己同型群にない次数スカラーが慣性として一つ余る。
一方、標識対象自身のスカラーをさらに慣性へ加える必要はない。
\end{lemma}

\begin{proof}
自然自己同型 $\eta$ の各 $t^{-n}\mathcal O$ 上の成分は、例外性により
スカラー $\lambda_n$ である。
$\Hom(t^{-n}\mathcal O,\mathcal O)=R_n\ne0$ と自然性から
$\lambda_n=\lambda_0$ となる。
反 Serre 軌道の豊富性により全ての対象はこれらの有限直和の商なので、
$\eta$ は全ての対象で同じスカラーである。
pointed 二射は標識上で恒等であるため、このスカラーは一に限られる。
特に $(\operatorname{Id},\lambda\operatorname{id}_{\mathcal O})$ と
$(\operatorname{Id},\operatorname{id}_{\mathcal O})$ は互いに二同型である。

graded 自己同型は AZ 再構成を通じて pointed exact 自己同値を与える。
逆に pointed exact 自己同値は Serre 関手と可換するので、
その可換同型を選ぶことで $R$ の graded 自己同型を与える。
可換同型の異なる選択は、前段の自然自己同型の計算により一つのスカラーだけ異なり、
$R_n$ 上ではその $n$ 乗だけ異なる。これで全射性を得る。

核を計算する。graded 自己同型 $\phi$ の誘導する pointed 関手が恒等と二同型なら、
各自由シフト $\pi R(i)$ 上のその二同型は、
基準同型 $b\mapsto\phi^{-1}(b)$ のスカラー $c_i$ 倍である。
隣接シフト間の Hom に対する自然性は
$c_{i+1}\phi_1=c_i\operatorname{id}_{R_1}$ を与える。
従って $\phi_1$ はスカラーで、$R$ の一次生成性により $\phi=\xi_\lambda$ である。
逆に次数スカラーでの加群 twist は、次数 $n$ 上で $\lambda^{-n}$ を掛ける写像により
恒等関手と自然同型であり、次数零の標識も保つ。
\end{proof}

九次元空間 $U$ を固定し、
\[
 \operatorname{Gr}\bigl(56,U\otimes U\bigr)
\]
上の関係部分束 $J$ から二次環 $R_J=T(U)/(J)$ を作る。
反 Serre 環として実現される部分を記号的に $X_{\mathrm{as}}$ と書く。
ここで $X_{\mathrm{as}}$ はまだ単なる幾何点集合以上の意味には構成していない。
$\GL(U)$ の中心は Grassmannian に自明に作用するため、
実現部分が有限表示の不変部分スキームとして構成できれば
\[
 [X_{\mathrm{as}}/\GL(U)]\longrightarrow
 [X_{\mathrm{as}}/\PGL(U)]
\]
は次数スカラーを除く $\mathbb G_m$-rigidification となる。
補題の幾何点慣性の計算から、pointed 圏のモジュライに対応すべきものは右辺である。
ただしこの式だけでは相対 $\HH$ モジュライの構成を完了しない。

\begin{proposition}[九生成元・二十五次元二次成分の gerbe は中立]
\label{prop:antiserre-gerbe-neutral}
$Y=\operatorname{Gr}(56,U\otimes U)$ とし、
$\mathcal J\subset U\otimes U$ を普遍部分束とする。
$[Y/\GL(U)]$ 上の普遍環の初めの二成分を
\[
 \mathcal R_1=U,\qquad
 \mathcal R_2=(U\otimes U)/\mathcal J
\]
と書く。その階数はそれぞれ九、二十五である。
線束
\[
 \mathcal L=(\det\mathcal R_1)^{-11}
                 \otimes(\det\mathcal R_2)^2
\]
は中心 $\mathbb G_m\subset\GL(U)$ に関して重み一である。
従って $\mathbb G_m$-gerbe
\[
 [Y/\GL_9]\longrightarrow[Y/\PGL_9]
\]
は中立であり、選んだ $\mathcal L$ によって
$[Y/\GL_9]\simeq[Y/\PGL_9]\times B\mathbb G_m$ と同値になる。
$\mathcal R_n\otimes\mathcal L^{-n}$ は $[Y/\PGL_9]$ 上へ降下し、
その全体は通常の準連接次数付き代数をなす。
特にこの商に残る独立の Brauer 障害はない。
同じ結論は任意の $\PGL_9$-不変部分への制限に成り立つ。
\end{proposition}

\begin{proof}
中心の元 $\lambda I$ は $\mathcal R_1$ に $\lambda$、
$\mathcal R_2$ に $\lambda^2$ で作用する。
従って $\det\mathcal R_1$ と $\det\mathcal R_2$ の重みは九と五十であり、
\[
 -11\cdot9+2\cdot50=1
\]
から $\mathcal L$ の重みは一である。
$\mathcal V=\mathcal R_1\otimes\mathcal L^{-1}$ の中心作用は自明なので、
$\GL_9$-同変降下により $[Y/\PGL_9]$ 上の階数九のベクトル束となる。
その射影枠束は普遍 $\PGL_9$-torsor と一致するので、
通常の枠束はこの torsor の $\GL_9$ への持ち上げを与える。
これが gerbe の中立化である。

同じ事実を局所的に見ると、$\PGL_9$-遷移関数の持ち上げ
$g_{ij}$ による三重重なりの障害が $\lambda_{ijk}$ であるとき、
$\mathcal L$ の遷移係数 $\ell_{ij}$ は重み一のため
$\ell_{ij}\ell_{jk}\ell_{ki}=\lambda_{ijk}$ を満たす。
従って $\ell_{ij}^{-1}g_{ij}$ は真の $\GL_9$-cocycle となる。
一般にこのような重み一の線束を持つ $\mathbb G_m$-gerbe は、
線束の枠束を取ることにより底の $B\mathbb G_m$ との積に同値となる。

$\mathcal R_n$ に対する中心の重みは $n$ なので
$\mathcal R_n\otimes\mathcal L^{-n}$ の重みは零である。
積もこの twist と両立する。
具体的には $\mathcal J\otimes\mathcal L^{-2}\subset
\mathcal V\otimes\mathcal V$ が部分束として降下し、
\[
 T(\mathcal V)/(\mathcal J\otimes\mathcal L^{-2})
\]
が求める代数である。
三次以上の成分の平坦性は Grassmannian 全体では主張せず、
別に定める相対 Hilbert 条件の locus で課す。
\end{proof}

\begin{remark}[有限型 locus と相対再構成の不足]
まず、幾何点の実現集合やその被約スキーム像を取るだけでは、
非被約底上の変形を指定したことにならない。
相対 Hilbert 関数・full section algebra 条件、AZ 豊富性、相対 Serre 同定、
平方根の存在と二次 AS 正則被覆の条件を、
有限表示の部分関手として記述する必要がある。
体上で成立する有限回の同値判定も、全底上で一様な次数の
Isom 関手や降下の表現可能性をそのまま与えるものではない。

直前の命題により、$\PGL_9$-torsor の Brauer 障害そのものは解消される。
しかし線束で次数ごとに twist した環の shift は、元の shift に
底からの線束をテンソルしたものに対応する。
従って相対 tails 圏の構成、相対 Serre 関手との同定、
その $\Hom(\mathcal O,t^n\mathcal O)$ の底変更を確認する仕事はなお残る。
gerbe の中立性は、この相対再構成を代用する主張ではない。

従って現段階で証明した内容は、有限 foundation の明示商スタック、
その Hochschild 計算、幾何点上の内在的二択、および反 Serre 表示で除くべき慣性の同定である。
これらを全 $\HH$ の有限型 Artin スタックと同一視する相対定理は残る。
\end{remark}
